\section{Affine limits of chart profiles}\label{sec:profiles}

The stationary increment estimate applies to invariant probabilities, whereas
the quasi-isometry under consideration need not be typical for any such
probability.  We first obtain invariant probabilities by averaging over large
charts.  Stationary increments then constrain the slopes of limiting chart
profiles to convex hulls of finitely many vectors.  A separate volume
comparison controls the support exponent of the image of each horizontal
segment.  Together these inequalities will make every supported limiting
profile affine on the interior of its chart, with linear part $A\pi$ for
some $A\in\mathcal A$.  The label and the affine offset may still differ
between profiles; their agreement is the subject of the next two sections.

Throughout this section the bounds on the map pairs are fixed, as in
Lemma~\ref{lem:pair-sections}.  They, and the group $G$, determine all
uniform constants below.  We assume $W\ne0$.  Fix a Euclidean norm on $W$,
and set
\begin{equation}\label{eq:profile-polytope}
 \mathcal C=\{w\in W:\alpha_i(w)\le\|\alpha_i\|
                         \text{ for every }i\}.
\end{equation}
Here the norm of a form is its dual Euclidean norm.  The unit ball is
contained in $\mathcal C$.  The positive relation among all generator
weights and their spanning property imply that $\mathcal C$ is bounded:
its recession cone consists of the vectors on which all the weights are
nonpositive, and that cone is $\{0\}$.  In particular $\mathcal C$ is a
compact convex polytope with $0$ in its interior.  Moreover
\[
 s_i(\mathcal C)=\|\alpha_i\|>0.
\]
Indeed the contact point of the unit ball for $\alpha_i$ belongs to
$\mathcal C$ and attains the defining bound.  At this contact point every
form on a different oriented ray has a strict inequality, by equality in
Cauchy--Schwarz.  Thus each distinct oriented generator ray defines a
genuine facet, and that contact point lies in its relative interior.

Write $m$ for the nilpotence class of $N$, and put
\[
 s_- =\min_i s_i(\mathcal C),\qquad
 s_+ =\max_i s_i(\mathcal C).
\]
We will sample an $N$-box with template $\mathcal C$ and allow an
additional deterministic $N$-translation from a slightly larger box.
To keep these samples asymptotically invariant under fixed group moves,
the $D$-positions must stay inside a smaller height region.  The next
choice ensures that this interior margin survives conjugation by the
additional translation.
Fix once and for all a number $c>1$ such that $cs_->2ms_+$.
Until the end of this section, $\delta$ is fixed and satisfies
$0<c\delta<1/2$.  Define the open subset of the graded nilpotent group
$D_\infty$ by
\begin{equation}\label{eq:profile-domain}
 \mathcal U_\delta
   =\{z\in D_\infty:\pi z\in(1-c\delta)\operatorname{int}(\mathcal C)\}.
\end{equation}
The map $\pi$ on $D_\infty$ is the linear height projection in graded
logarithmic coordinates.  Consequently $\mathcal U_\delta$ is connected.

\begin{definition}[Chart profiles]\label{def:chart-profiles}
Let $r\to\infty$ along a sequence.  For each $r$, choose a map pair with
first map $F$, a basepoint $x\in G$, and a deterministic suffix $k_0\in N$
satisfying
\begin{equation}\label{eq:profile-suffix}
 k_0\in B_r((1+\delta)\mathcal C,e^{\varepsilon_r r}),
 \qquad \varepsilon_r\ge0,\quad \varepsilon_r\longrightarrow0.
\end{equation}
Choose $a$ uniformly with respect to Haar measure on $B_r(\mathcal C)$,
and set $n=ak_0$.  The chart profile is the random continuous function
\begin{equation}\label{eq:chart-profile}
 f_{r,x,n}(z)=\frac{\pi F(xn\Delta_r z)-\pi F(x)}{r},
 \qquad z\in D_\infty.
\end{equation}
The law of this function is taken in
$C_{\mathrm{loc}}(D_\infty,W)$, with its topology of uniform convergence
on compact subsets.  The map, basepoint, and suffix may all vary with $r$.
\end{definition}

For a continuous $W$-valued function $h$ on a nonempty compact set $Z$,
write $\|h\|_Z=\sup_{z\in Z}\|h(z)\|$.

The logarithmic coordinates of $n$ are bounded by $e^{O(r)}$, uniformly
along every sequence in Definition~\ref{def:chart-profiles}.  Hence
Lemma~\ref{lem:log-distortion} gives $|n|_G=O(r)$, so the values of the
profiles at the identity are bounded.  Lemma~\ref{lem:d-scaling} and the
uniform Lipschitz bound for $F$ give equicontinuity and boundedness on
each compact subset of $D_\infty$: explicitly,
\[
 \|f_{r,x,n}(z)-f_{r,x,n}(z')\|
       \le \frac{C_1}{r}\,d_D(\Delta_r z,\Delta_r z'),
\]
where $C_1$ is independent of $a,x,F,k_0$.  Arzel\`a--Ascoli and a compact
exhaustion therefore give a common compact set containing the profiles
along any such sequence, after discarding finitely many terms.
In particular their laws have weakly convergent subsequences.  Every
limiting function is Lipschitz along horizontal segments, with a bound
proportional to the norm of the horizontal velocity.  Indeed those
segments are the uniform limits of
\begin{equation}\label{eq:actual-profile-segment}
 \Psi_{r,t}(z)=\Delta_r^{-1}(\Delta_r z\exp_D(rt v)),
 \qquad v\in V,
\end{equation}
and the corresponding paths in $D$ have speed at most a constant times
$r\|v\|$.  Here, as in Lemma~\ref{lem:d-scaling}, a vector in $V$ is
represented in the chosen first-layer complement of $\mathfrak d$.

\subsection{Invariant occupation probabilities}

We need chart averages that are asymptotically unchanged by each fixed
right translation.  The enlarged suffix in
\eqref{eq:profile-suffix} is important for later comparisons, so the
following estimate keeps track of its effect.

\begin{lemma}[Interior slack]\label{lem:profile-interior-slack}
Use the data of Definition~\ref{def:chart-profiles}.  Suppose a family
$u_r\in N$ has a word expansion of its logarithm whose coefficient of
each word $Z$ is bounded in absolute value by
\[
 \exp\bigl(r(1-c\delta)h_Z(\mathcal C)+o(r)\bigr),
\]
with a common error envelope for the finitely many words.  Then
$k_0u_rk_0^{-1}$ has logarithmic coordinates tending exponentially to
zero in the coordinates normalized for $B_r(\mathcal C)$.
The same conclusion holds uniformly for bounded families satisfying
the displayed coefficient bounds with a common envelope.
\end{lemma}

\begin{proof}
Expand $\log k_0$ in the finite list of bracket words using
Lemma~\ref{lem:box-calculus}.  Its coefficient in a word $Y$ is bounded
by $\exp(r(1+\delta)h_Y(\mathcal C)+o(r))$.  The finite expansion of
$\exp(\ad(\log k_0))\log u_r$ consists of brackets of an initial word
$Z$ with additional words $Y_1,\ldots,Y_j$.  The sum of their lengths
is at most $m$ whenever the bracket is nonzero.  Before expansion in a
maximizing basis, the logarithm of the coefficient divided by $r$ is
at most
\[
 (1-c\delta)h_Z(\mathcal C)
       +(1+\delta)\sum_{\ell=1}^j h_{Y_\ell}(\mathcal C)+o(1).
\]
The width of the resulting word is
$h_Z(\mathcal C)+\sum_\ell h_{Y_\ell}(\mathcal C)$.  Relative to this
width the exponent is therefore at most
\begin{equation}\label{eq:profile-slack-exponent}
 -\delta\left(c h_Z(\mathcal C)
                  -\sum_{\ell=1}^j h_{Y_\ell}(\mathcal C)\right)+o(1)
 \le -\delta(cs_--ms_+)+o(1).
\end{equation}
This is bounded above by a negative constant for all large $r$.
The basis-exchange comparison in Lemma~\ref{lem:box-calculus} says
that expanding a word in a maximizing basis uses only basis entries
whose widths are at least its width.  That expansion cannot increase
the bound in \eqref{eq:profile-slack-exponent}.  There are only finitely
many terms, which proves the assertion.
\end{proof}

For a pair $\omega$ and a point $p\in G$, let $\omega[p]\in X$ denote
the pair normalized at $p$; its first map is
$g\mapsto F(p)^{-1}F(pg)$.  The normalized action has the exact
identity
\begin{equation}\label{eq:profile-normalization-action}
 T_g\omega[p]=\omega[pg].
\end{equation}

\begin{lemma}[Occupation invariance]\label{lem:occupation-invariance}
Let $\rho$ be a continuous compactly supported probability density
with respect to Haar measure on $D_\infty$, with support contained in
$\mathcal U_\delta$.  Independently of $a$, sample $z$ with density
$\rho$.  Every weak limit of the laws of
\[
 \omega[xak_0\Delta_r z]\in X
\]
is an invariant probability for the normalized right action of $G$.
\end{lemma}

\begin{proof}
We prove that the laws of the sampled points are asymptotically
unchanged in total variation by right multiplication by each fixed
element of $G$.  It suffices to consider elements of $D$ and $N$,
since these subgroups generate $G$.

For $d\in D$, right multiplication changes $z$ to
\[
 z\longmapsto\Delta_r^{-1}(\Delta_r z\,d).
\]
This map preserves Haar measure in the scaled coordinates, because
$D$ is unimodular and the coordinate scaling has constant Jacobian.
By Lemma~\ref{lem:d-scaling} it converges to the identity uniformly
on compact sets.  A continuous compactly supported density changes
by $o(1)$ in $L^1$, proving the assertion for $d$.

For fixed $u\in N$, right multiplication changes $a$ to
\[
 a\,q_r(z),\qquad
 q_r(z)=k_0\Ad_{\Delta_r z}(u)k_0^{-1},
\]
and leaves $z$ unchanged.  On the compact support of $\rho$, the
weight-action estimate in Lemma~\ref{lem:canonical-height} bounds a
coefficient of $\log\Ad_{\Delta_r z}(u)$ in a word $Z$ by
\[
 \exp\bigl(r\gamma_Z(\pi z)+o(r)\bigr)
 \le \exp\bigl(r(1-c\delta)h_Z(\mathcal C)+o(r)\bigr).
\]
The error is uniform there; it includes the polynomial factors from
the unipotent part of the $D$ action.  Words of zero total weight
cause no exception to the inequality.  Lemma~\ref{lem:profile-interior-slack}
shows that $q_r(z)$ tends exponentially to the identity in the
$\mathcal C$-normalized coordinates, uniformly in $z$.

The normalized product laws have uniformly bounded polynomial
coefficients.  Thus, on the fixed coordinate cube defining
$B_r(\mathcal C)$, multiplication by $q_r(z)$ and by $q_r(z)^{-1}$
both converge uniformly to the identity.  Their images lie in
Euclidean boundary enlargements of the cube of thickness tending
to zero.  Haar invariance now gives
\[
 \sup_{z\in\supp\rho}
 \frac{\vol_N(B_r(\mathcal C)q_r(z)
                     \mathbin{\triangle} B_r(\mathcal C))}
      {\vol_N(B_r(\mathcal C))}\longrightarrow0.
\]
This is precisely the required total-variation estimate for the
conditional uniform laws of $a$.

The estimates were proved before projecting the sampling coordinates
to $G$.  Projection decreases total variation, so the possible
nonuniqueness of $ND$ coordinates is harmless here.  Finally apply
\eqref{eq:profile-normalization-action}.  The action on the compact
space $X$ is continuous by Lemma~\ref{lem:pair-sections}; every weak
limit of the normalized occupation laws is therefore invariant.
\end{proof}

\subsection{Two inequalities on a limiting segment}

Fix a supported limit $f$ of chart profiles.  We first constrain
its increments to the convex hull of the stationary slopes, then
compare the support exponent of each image segment with that of its
source.  Finite derivative labels alone would still permit folding:
the scalar function $t\mapsto|t|$ has derivative $\pm1$ almost everywhere.
The second inequality retains the volume information needed to rule
out this behavior.  Lemma~\ref{lem:affinity-criterion} will combine
the two inequalities to prove affinity.
Both statements concern every compact horizontal segment in the
open chart, and neither uses any comparison between neighboring
sheets.

\begin{lemma}[Convex increment constraint]\label{lem:profile-convex-increments}
Let $\nu$ be a weak limit of chart-profile laws, and let
$f\in\supp\nu$.  If
\[
 J=\{z\exp_\infty(tv):0\le t\le T\}\Subset\mathcal U_\delta,
 \qquad T>0,\quad v\in V,
\]
then
\begin{equation}\label{eq:profile-convex-increments}
 \frac{f(z\exp_\infty(Tv))-f(z)}{T}
       \in\operatorname{conv}\{A\pi v:A\in\mathcal A\}.
\end{equation}
\end{lemma}

\begin{proof}
First allow the starting point $z$ to vary in a small neighborhood
$O$ so that the whole family of segments lies in a compact subset
of $\mathcal U_\delta$.  Choose a continuous probability density
$\eta$ supported in $O$, positive near the specified starting point.
Choose another continuous probability density $\rho$, compactly
supported in $\mathcal U_\delta$, which is bounded below by a
positive number on a compact neighborhood of this swept family.
The same neighborhood contains all the actual scaled paths
$\Psi_{r,t}(z)$ from \eqref{eq:actual-profile-segment}, for all large
$r$, $z\in\supp\eta$, and $0\le t\le T$.

Let $\mu_r$ be the occupation law obtained with $\rho$.  Pass to a
subsequence on which $\mu_r$ converges to $\mu$; the prescribed
profile limit remains $\nu$.  By Lemma~\ref{lem:occupation-invariance},
$\mu$ is invariant.  The cutoff and this subsequence are fixed
before choosing a microstep length.

Fix an integer $L\ge1$.  Put $q=\lfloor rT/L\rfloor$ and divide the
actual path into $q$ increments by $\exp_D(Lv)$ and a final remainder.
For $0\le j<q$, its $j$th starting point is
\[
 p_j=xak_0\Delta_r z\exp_D(jLv).
\]
Right invariance of Haar measure in $D$ shows that the distribution
of the scaled $D$ coordinate of $p_j$, when $z$ has density $\eta$,
has density at most $\|\eta\|_\infty$.  Its support is in the
compact set where $\rho$ has a positive lower bound.  Hence the
law of $\omega[p_j]$ is bounded above by $C_0\mu_r$, with a constant
$C_0$ independent of $r,j,L$.  The original law of $a$ is unchanged
and independent of $z$ throughout this comparison.

Write $K_v=\operatorname{conv}\{A\pi v:A\in\mathcal A\}$ and define
the bounded continuous function on $X$
\[
 d_L(\omega')=
 \dist\left(\frac{\pi F_{\omega'}(\exp_D(Lv))}{L},K_v\right).
\]
The identity of height cocycles at $p_j$ and convexity give
\begin{align*}
 &\mathbb E_{a,z}\dist\left(
 \frac{f_{r,x,ak_0}(\Psi_{r,T}(z))-f_{r,x,ak_0}(z)}{T},K_v\right)
 \\
 &\hspace{25mm}\le
   \frac{L}{rT}\sum_{j=0}^{q-1}\mathbb E_{a,z}d_L(\omega[p_j])
       +O(L/r)
 \le C_0\int_X d_L\,d\mu_r+O(L/r).
\end{align*}
The final error uses the Lipschitz bound on the remainder interval
and the boundedness of $K_v$.  For this fixed $L$, let $r\to\infty$.
Profile compactness and the uniform convergence of $\Psi_{r,T}$ give
\[
 \int\!\int_O
 \dist\left(\frac{f'(z\exp_\infty(Tv))-f'(z)}{T},K_v\right)
      \eta(z)\,dz\,d\nu(f')
 \le C_0\int_X d_L\,d\mu.
\]
Now let the integer $L$ tend to infinity.
Corollary~\ref{cor:stationary-increments} for this one fixed invariant
probability $\mu$ makes the right side tend to zero.  No convergence
rate uniform over invariant probabilities is required.

The nonnegative integrand on the left is continuous in the profile
and the starting point.  Its integral being zero, it vanishes on
the product of the supports.  In particular it vanishes for the
specified $f$ and $z$.  The argument applies to each segment as
stated in the lemma.
\end{proof}

\begin{lemma}[Segment volume inequality]\label{lem:profile-segment-volume}
Under the hypotheses of Lemma~\ref{lem:profile-convex-increments},
\begin{equation}\label{eq:profile-segment-volume}
 P(f(J))\ge P(\pi J)=T p(\pi v).
\end{equation}
\end{lemma}

\begin{proof}
\emph{Fixing the endpoint scheme.}
Set $E=\pi J$.  We may pass to any further subsequence of the one
giving $\nu$.  First pass to a subsequence on which the nilpotent
laws in $E$-normalized coordinates converge.  By
Lemma~\ref{lem:product-schemes}, choose a fixed finite ordered list
of generator increments whose normalized endpoint map has a
submersive limit at an interior parameter point.  Its length $h$
is fixed independently of $r$.  An increment in letter $i$ has the
form $\exp(t e^{rs_i(E)}X_i)$, with $t$ in a fixed bounded interval.
Replace $t$ by $t_1-t_2$ with independent uniform interval parameters,
enlarging their fixed intervals if necessary.  The difference map
is surjective, so the limiting endpoint map is still submersive
somewhere in the parameter cube.  Fix this scheme now.

\emph{Retaining sheets with the prescribed profile.}
Choose a compact neighborhood $Z_J\Subset\mathcal U_\delta$ of $J$
that contains all the actual scaled paths $\Psi_{r,t}(z)$,
$0\le t\le T$, for large $r$.  Put $Q_r=B_r(\mathcal C)k_0$.
Because $f\in\supp\nu$, for
each fixed positive tolerance the set of $n\in Q_r$ whose profiles
are within that tolerance of $f$ on this compact has a positive
lower limiting relative measure.  Choosing the tolerances slowly
enough gives Borel sets $S_r\subset Q_r$ and numbers $\zeta_r\to0$
such that
\begin{equation}\label{eq:profile-good-mass}
 \frac{\vol_N(S_r)}{\vol_N(Q_r)}\ge e^{-o(r)},\qquad
 \sup_{n\in S_r}\|f_{r,x,n}-f\|_{Z_J}\le\zeta_r.
\end{equation}
To be explicit, choose tolerances $1/j$ successively; postpone the
$j$th choice until its probability lower bound has logarithm of
absolute value at most $r/j$.  This also ensures that every other
fixed compactness error tends to zero along the diagonal.  The
scheme and its compact were chosen before this diagonal.

We verify the support-volume hypothesis for recurrence.  For an
individual generator interval let $\sigma_{j,r}$ denote the law
of $\exp(t e^{rs_i(E)}X_i)$.  Since
$E\subset(1-c\delta)\operatorname{int}(\mathcal C)$,
\[
 s_i(E)\le(1-c\delta)s_i(\mathcal C).
\]
Lemma~\ref{lem:profile-interior-slack} shows that
$k_0\supp(\sigma_{j,r})k_0^{-1}$ tends to the identity in normalized
$\mathcal C$ coordinates.  The bounded polynomial product law
therefore gives
\begin{equation}\label{eq:profile-recurrence-denominator}
 \vol_N(Q_r\supp(\sigma_{j,r}))
       \le D_{j,r}\vol_N(Q_r),\qquad D_{j,r}=e^{o(r)}.
\end{equation}
All estimates hold with common envelopes for the fixed finite list.

For the $j$th switch, sample $s_{j,1},s_{j,2}$ independently with law
$\sigma_{j,r}$, independently also for different $j$, and put
$u_j=s_{j,1}s_{j,2}^{-1}$.  Its law is
$\sigma_{j,r}*\check\sigma_{j,r}$.
Apply Lemma~\ref{lem:backward-recurrence} in the unimodular group $N$
with the compact set $Q_r$, the Borel set $S_r$, and these increment
laws.  The mass bound \eqref{eq:profile-good-mass}, the denominator
bound \eqref{eq:profile-recurrence-denominator}, and the fixed depth
$h$ give a starting point $n_0\in S_r$ such that
\[
 \mathbb P\{n_0u_1\cdots u_j\in S_r
                   \text{ for every }1\le j\le h\}\ge e^{-o(r)}.
\]
Thus every intermediate sheet in a successful sequence has a good
profile.

\emph{From parameter measure to endpoint volume.}
The successful interval parameters have
Lebesgue measure at least $e^{-o(r)}$ in the fixed product cube,
since the uniform product density is a fixed constant.  In
$E$-normalized logarithmic coordinates, the endpoint map for the
increments is a polynomial of fixed degree with convergent bounded
coefficients and a submersive limit, by our choice of scheme.
Lemma~\ref{lem:polynomial-image}, or equivalently
Corollary~\ref{cor:restricted-products}, now gives accessible
endpoint volume in $N$ at least
\begin{equation}\label{eq:profile-accessible-volume}
 \exp\bigl(rP(E)-o(r)\bigr).
\end{equation}
If necessary, restrict the successful parameters to a compact subset
of half their measure; this preserves the estimates and makes the
image measurable.  Left multiplication by $n_0$ preserves the
volume.  The rank conclusion is a property of the full fixed
polynomial map, established before the successful restriction.

\emph{Comparing source and image packing.}
Take $d_p=\Delta_r z$ as pivot.  For each prescribed generator increment,
travel along $\Delta_r z\exp_D(rt v)$ to a point where its weight
attains $s_i(E)$, perform the increment, and return to the pivot on
the new sheet.  Lemma~\ref{lem:support-switch-paths} gives an
$o(r)$-length switch at that point; the two $D$ legs have length
$O(r)$.  The difference of two interval parameters is a single
generator increment.  We therefore do not need a good profile at
the algebraic midpoint of $s_1s_2^{-1}$.

All sheet states at which a $D$ leg is traversed belong to $S_r$.
The images of these legs consequently have heights in
\[
 \pi F(x)+r\bigl(f(J)+o(1)\bigr).
\]
The short switches contribute only $o(r)$ height errors by
Lipschitzness.  The finite concatenation has length $O(r)$, begins
at $xn_0d_p$, and ends at the fixed $D$ position $d_p$ on one of the
accessible sheets.

By Lemma~\ref{lem:fixed-sheet-packing}, the endpoint volume
\eqref{eq:profile-accessible-volume} gives at least
$\exp(rP(E)-o(r))$ source points separated by any prescribed fixed
large constant.  Choose this constant so their quasi-isometry
images remain separated.  Their image paths have a common start,
length $O(r)$, and relative heights in
$r(f(J)-f(z)+o(1))$, since the starting profile value is
$f_{r,x,n_0}(z)=f(z)+o(1)$.  The compact set $f(J)-f(z)$ contains
zero; enlarging it by a ball of radius tending to zero absorbs the
error.  Lemma~\ref{lem:path-packing}, translation invariance
of $P$, and its continuity for Hausdorff distance give the upper
bound $\exp(rP(f(J))+o(r))$ for the same separated points.
Taking logarithms and dividing by $r$ proves
\eqref{eq:profile-segment-volume}.  When $\pi v=0$ the claimed lower
bound is zero; it is also immediate from nonnegativity of $P$.
\end{proof}

\subsection{A deterministic affinity criterion}

We now use the two segment inequalities twice.  At infinitesimal
scale they force the horizontal derivative to be $A\pi$ for a
measurable label $A\in\mathcal A$.  On whole segments, equality
between support and positive variation fixes the weight signs along
directions avoiding every weight hyperplane.  The closures of the cones
determined by these signs have extremal rays; passage to those directions
gives enough affine line restrictions to obtain smoothness.  The derivative label
is consequently continuous, hence constant on the connected domain.
The following criterion isolates this deterministic argument.

\begin{lemma}[Affinity criterion]\label{lem:affinity-criterion}
Let $f:\mathcal U_\delta\to W$ be continuous and locally Lipschitz
along horizontal segments, with locally uniform bounds for unit
horizontal velocities.  Suppose that, for every segment
$J=\{z\exp_\infty(tv):0\le t\le T\}\Subset\mathcal U_\delta$,
$T>0$, both
\begin{equation}\label{eq:affinity-two-inequalities}
 \frac{f(z\exp_\infty(Tv))-f(z)}{T}
       \in\operatorname{conv}(\mathcal A\pi v),
 \qquad P(f(J))\ge T p(\pi v)
\end{equation}
hold.  Then there exist $A\in\mathcal A$ and $b\in W$ such that
$f(z)=A\pi z+b$ for all $z\in\mathcal U_\delta$.
\end{lemma}

\begin{proof}
\emph{Horizontal derivatives take values in the finite group.}
Choose a basis of the first layer $V$ and its left-invariant vector
fields on $D_\infty$.  Flow-box Fubini and linewise Lipschitzness
give locally bounded distributional derivatives of $f$ in these
directions.  Their linear combinations define a measurable linear
map $M(z):V\to W$.  For each fixed $v$, it agrees almost everywhere
with the line derivative in direction $v$: both are the
distributional derivative along that field.  These assertions may
be made simultaneously for a countable dense set of $v$ by
intersecting their full Haar-measure sets.

At a point of line differentiability, the image of a short segment,
translated to start at zero and divided by its length, converges
in Hausdorff distance to the linear segment of velocity $M(z)v$.
Indeed the definition of differentiability bounds the error
uniformly on the entire sufficiently short one-sided interval.
Differentiate \eqref{eq:affinity-two-inequalities} and use continuity
and positive homogeneity of $P$ to obtain
\begin{equation}\label{eq:affinity-derivative-constraints}
 Mv\in\operatorname{conv}(\mathcal A\pi v),\qquad
 p(Mv)\ge p(\pi v)
\end{equation}
at the common full-measure points.  Continuity in $v$ extends these
constraints from the dense set to all of $V$.

The first constraint makes $M$ vanish on $\ker\pi$, so $M=B\pi$
for a linear map $B:W\to W$.  It also gives
$p(Bw)\le p(w)$ by convexity, while the second constraint gives the
reverse inequality.  Thus $B$ is a linear isometry of $p$ and is
invertible.  The full linear isometry group of the polyhedral norm
$p$ is finite: its elements permute the finitely many vertices of
the unit ball, and those vertices span $W$.  Average a Euclidean
inner product over this full finite group.  In its norm, $Bw$ and
every $Aw$, $A\in\mathcal A$, have the same length.  Strict
convexity and the convex-hull inclusion force
\[
 Bw\in\{Aw:A\in\mathcal A\}\qquad(w\in W).
\]
If $B$ differed from every $A\in\mathcal A$, the finitely many
proper linear subspaces $\ker(B-A)$ could not cover $W$.  Hence
$B\in\mathcal A$.  We have proved
\begin{equation}\label{eq:affinity-label-field}
 M(z)=B(z)\pi,\qquad B(z)\in\mathcal A
 \quad\text{for almost every }z.
\end{equation}
Choose any label on the remaining null set, so $B$ is a measurable
finite-valued function throughout the domain.

\emph{Equality of support and variation fixes weight signs.}
Let $v$ have generic height $w=\pi v$, meaning
$\alpha_i(w)\ne0$ for all $i$.  Along almost every line of this
direction, the absolutely continuous restriction
$h(t)=f(z\exp_\infty(tv))$ has derivative $B(z\exp_\infty(tv))w$
almost everywhere.  Thus $p(h'(t))=p(w)$.  We work on a compact
interval $[0,T]$ inside such a line and translate $h$ to have
$h(0)=0$, which does not change $P$.

For a word basis $\mathcal B$, let $m_i(\mathcal B)$ count all
occurrences of $X_i$ in its words.  Every generator is mandatory,
so $m_i(\mathcal B)\ge1$.  Choose a basis maximizing
$P(h([0,T]))$, and put $q_i(t)=(\alpha_i\circ h)'(t)$.  Then
\begin{align}
 P(h([0,T]))
 &=\sum_i m_i(\mathcal B)\max_{0\le t\le T}\alpha_i(h(t))
 \nonumber\\
 &\le\sum_i m_i(\mathcal B)\int_0^T(q_i(t))_+\,dt
 \le\int_0^T p(h'(t))\,dt
 =T p(w).
 \label{eq:affinity-support-variation}
\end{align}
The second inequality is the definition of $p$ as the maximum over
word bases.  The segment volume inequality supplies the reverse
inequality between the outer quantities, so every inequality in
\eqref{eq:affinity-support-variation} is equality.  Since every
$m_i(\mathcal B)$ is positive, each scalar maximum has no deficit.
If $t_i$ is a maximizing time for $\alpha_i(h(t))$, this says
\begin{equation}\label{eq:affinity-sign-order}
 0=\int_0^{t_i}(q_i(t))_-\,dt
                  +\int_{t_i}^T(q_i(t))_+\,dt.
\end{equation}
Consequently a negative derivative cannot occur, on a set of
positive measure, before a positive derivative.  On this interval
each weight's positive-sign indicator is nonincreasing, up to null
sets.

By Lemma~\ref{lem:support-symmetries}, pullback by every label
permutes the distinct oriented generator-weight rays.  Genericity
therefore makes every $q_i$ nonzero almost everywhere, and the
number of positive signs among those distinct rays is independent
of the label $B$.  A finite collection of nonincreasing indicators
with constant sum has each indicator constant almost everywhere.
Thus all weight signs are constant along the line, on the interval
under consideration.  Applying this on a countable exhaustion of
its open interval in $\mathcal U_\delta$ shows the following:
if $w$ lies in a chamber of the generator-weight hyperplane
arrangement, then almost all labels along the line send that
chamber to the same chamber.  Here each line meets
$\mathcal U_\delta$ in an interval because its projected height
moves linearly through a convex set.

\emph{Extremal directions have constant horizontal derivative.}
The chambers just used are the connected components of
$W\setminus\bigcup_i\ker\alpha_i$.  Their closures are pointed
polyhedral cones, since the weights span $W^*$.  Let $w\ne0$ span
an extremal ray $R$ of one of these closed cones, and choose a
horizontal lift $v\in V$.

For each chamber $C_0$ whose closure contains $R$, choose generic
vectors $v_j\to v$ with $\pi v_j\in C_0$.  We explain carefully
why the almost-everywhere chamber agreement passes to this limiting
direction.  On a compact set of pairs $(z,s)$ whose flows stay in
the domain, the maps $z\mapsto z\exp_\infty(sv_j)$ preserve Haar
measure and converge uniformly to the corresponding map for $v$.
Approximating the bounded measurable function $B$ by continuous
functions on a larger compact set shows that
\[
 B(z\exp_\infty(sv_j))
       \longrightarrow B(z\exp_\infty(sv))
 \quad\text{in measure locally in }(z,s).
\]
The approximation errors are uniformly controlled by Haar
invariance.  Apply the same statement at a second time $t$;
Fubini gives convergence in measure locally in $(z,s,t)$.
After a subsequence both labels converge almost everywhere.
We restrict to compact sets of triples for which the segment
between the two limiting positions stays in the open domain;
for large $j$ the nearby segments do so too.  The preceding
chamber argument applies to every $v_j$.  Since the set of labels
is finite, the relation that two labels have the same image of
$C_0$ is closed.  It follows that almost every such triple satisfies
\[
 B(z\exp_\infty(sv))\overline C_0
   =B(z\exp_\infty(tv))\overline C_0.
\]
There are only finitely many incident chambers.  Intersect their
full-measure sets, and then exhaust the domain of triples by
compact sets.

The intersection of the closures of all chambers incident to $R$
is exactly $R$.  Indeed the active hyperplanes of an extremal ray
cut down to its line; the remaining chamber inequalities select
the half-line containing $w$.  Equivalently, all hyperplanes
through $w$ allow both neighboring signs among the incident
chambers, while the other hyperplanes have the fixed sign of $w$.
Intersecting the displayed equalities therefore shows that the
two labels have the same image of $R$.  They preserve the invariant
Euclidean norm chosen above, so their values on $w$ are equal.
Thus the derivative in direction $v$ is constant at almost every
pair of positions on almost every such line.  Absolute continuity
makes $f$ affine on those lines.

When $v\in\ker\pi$, the first inequality in
\eqref{eq:affinity-two-inequalities} already makes $f$ constant
on every horizontal segment of direction $v$.

\emph{Regularity makes the label constant on the domain.}
The extremal rays of a full-dimensional pointed chamber span $W$.
Choose finitely many lifts of such rays spanning a complement to
$\ker\pi$ in $V$, and add a basis of $\ker\pi$.  The resulting
vectors $v_1,\ldots,v_b$ span $V$.  Let $Y_j$ be their
left-invariant horizontal fields on $D_\infty$.  The line
affinities just proved imply
\[
 Y_j^2 f=0\quad\text{in distributions},\qquad 1\le j\le b.
\]
For completeness, use flow coordinates for $Y_j$ and test against
a compactly supported smooth function.  Along almost every flow
line the second distributional derivative of an affine function
is zero.  Integration in the transverse coordinates gives the
displayed identity.  The flow preserves Haar measure, so no
divergence term occurs.

The associated graded Lie algebra of $D$ is generated by its first
layer $V$.  Hence the fields $Y_j$ and their iterated brackets span
the full tangent space everywhere.  H\"ormander's sum-of-squares
theorem \cite[Theorem~1.1]{Hormander1967} applies to
$\sum_{j=1}^bY_j^2$, componentwise in $W$, and makes $f$ smooth.
All hypotheses are local on the open domain: the fields are smooth
and real, the bracket condition holds at every point, and the
right side is the smooth function zero.

The horizontal derivative is now continuous.  Its almost-everywhere
membership in the finite set $\{A\pi:A\in\mathcal A\}$ from
\eqref{eq:affinity-label-field} extends to every point, since a
full-measure subset is dense.  The domain is connected, so that
derivative is one constant $A\pi$.  Every horizontal derivative
of $f-A\pi$ vanishes; the same is then true for iterated brackets.
Bracket generation makes its entire differential zero, and
connectedness gives $f-A\pi=b$ for one constant $b\in W$.
\end{proof}

\begin{proposition}[Affine profile limits]\label{prop:affine-profiles}
For every fixed $\delta$ with $0<c\delta<1/2$ and every sequence
in Definition~\ref{def:chart-profiles}, the profile laws are
relatively compact.  Every function in the support of every
limiting law has the form
\begin{equation}\label{eq:affine-profile-limit}
 f(z)=A\pi z+b\quad(z\in\mathcal U_\delta),
 \qquad A\in\mathcal A,\quad b\in W.
\end{equation}
\end{proposition}

\begin{proof}
Relative compactness and horizontal Lipschitz regularity were
proved after Definition~\ref{def:chart-profiles}.
Lemmas~\ref{lem:profile-convex-increments} and
\ref{lem:profile-segment-volume} supply both assumptions of
Lemma~\ref{lem:affinity-criterion} on every segment.  That lemma
gives \eqref{eq:affine-profile-limit}.
\end{proof}

We record precisely the uniform consequence of this supported-law
statement that will be used for additional suffixes.  It does not
yet select a single affine function for a whole chart.

\begin{corollary}[Distance to the affine class]\label{cor:affine-profile-concentration}
Fix $\delta$ as above, a function $\varepsilon(r)\ge0$ tending to
zero, a nonempty compact set $Z\subset\mathcal U_\delta$, and
$\eta>0$.  Uniformly over the fixed family of map pairs, all
$x\in G$, and all deterministic
$k_0\in B_r((1+\delta)\mathcal C,e^{r\varepsilon(r)})$,
\begin{equation}\label{eq:distance-to-affine-class}
 \mathbb P_a\left\{
  \min_{A\in\mathcal A}\inf_{b\in W}
     \sup_{z\in Z}\|f_{r,x,ak_0}(z)-A\pi z-b\|\ge\eta
                 \right\}\longrightarrow0.
\end{equation}
The conclusion also holds when $k_0$ is sampled independently of
$a$ from any probability supported in the same allowed suffix
set.  In that case every supported limiting profile is still
affine on $\mathcal U_\delta$.
\end{corollary}

\begin{proof}
For a continuous function on $Z$, its distance in the supremum
norm to the set of restrictions of $A\pi+b$ is a continuous,
$1$-Lipschitz function of that continuous function.  The affine
restriction set is closed: take a subsequence with fixed $A$,
and evaluate at one point of $Z$ to see that the offsets converge.

If \eqref{eq:distance-to-affine-class} failed, there would be a
sequence of maps, basepoints, and deterministic allowed suffixes
for which the displayed probability is bounded below by a positive
constant.  Profile compactness supplies a subsequential law limit.
Proposition~\ref{prop:affine-profiles} supports that limit on the
affine class.  The event whose probability is bounded below is
closed, so the Portmanteau inequality is a contradiction.

For an independent random suffix, condition on its value.  The
bad probability is an average of deterministic bad probabilities,
and hence is bounded by the same supremum, which tends to zero.
For a compact exhaustion of $\mathcal U_\delta$, the zero sets of
the continuous affine-distance functions are closed.  For each
compact in the exhaustion and each $\eta>0$, the Portmanteau
inequality for the open event that this distance exceeds $\eta$
shows that its limiting probability is zero.  Thus every limiting
law is supported on functions affine on all these compact sets.  Use a
constant-label subsequence and evaluate at a fixed point to see
that their labels and offsets are compatible across the exhaustion;
equivalently, the class of functions affine on the entire domain
is closed in $C_{\mathrm{loc}}(\mathcal U_\delta,W)$.

The function $\varepsilon(r)$ was fixed before taking the limit.
This assertion is not a supremum over all possible subexponential
envelopes at one finite scale.  Nor does it assert concentration
near one particular pair $(A,b)$: different affine profiles may
still occur in the limiting law.
\end{proof}
